\documentclass{article}
\usepackage{hyperref}
\usepackage{amsmath}
\usepackage{amssymb}
\usepackage{graphicx}
\usepackage{parskip}
\usepackage{array}
\usepackage{booktabs}
\usepackage{verbatim}

\title{Formation, Yield, and Persistence --- Evolution by Emergence as an Optimization Principle}
\author{}
\date{}

\begin{document}

\maketitle

\subsection*{A working note: cost and barrier from nucleation, yield reconceived as the \textit{reshaping of reachability}, persistence as a self-maintaining loop --- with the gradient made local and non-teleological}

\begin{quote}
\textit{
Status: the framework is \textbf{not} "everything is nucleation." It is an \textbf{optimization principle} --- structures persist when what they release repays the cost of bringing and maintaining them into existence --- for which classical nucleation theory is the one \textbf{rigorously solved instance} (crystals). Building the theory on the principle and exhibiting nucleation as a realization of it is far more robust than building it on the analogy: the principle survives even if a mapping turns out imperfect. The claim is the careful one --- \textbf{many adaptive systems appear to exhibit nucleation-like barrier dynamics} --- not the universal one.

Two refinements reorganize the back half of this note. First, \textbf{yield is released slack}: the real, budgeted, conserved capacity a structure frees when one of its loops runs more cheaply (free energy, labour-time, capital, attention --- §10). Slack is a \textit{scalar}, and that is not a problem --- it is a \textit{present, local, conserved} quantity, not a "future value." The teleology trap was never the scalar-ness of the \textit{released} quantity; it was trying to scalarize the \textit{downstream consequence.} That consequence --- the \textbf{reshaping of which futures become reachable} --- is \textit{structural, not scalar}, and it is what the released slack \textit{does}, not what yield \textit{is} (§9). The corrected chain: cheaper loop $\rightarrow$ released slack (= yield) $\rightarrow$ barriers lowered $\rightarrow$ reachability reshaped $\rightarrow$ new loops possible.

Second, the mechanism is \textbf{slack-flow through loops}, and "flows toward" is formalized as \textit{myopic diffusion of a budgeted quantity} on the edge graph (§10), removing the last suspicion of teleology by construction.

Where this note states logical entailments ("if X enters the cost term, the cube law implies Y"), those are kept crisp; hedging a deduction as one hedges an observation blurs the line between what the math entails and what we have observed, and that line is what a research program must keep sharp.
}
\end{quote}

\begin{quote}
\textit{
On the status of this document: it is the dense \textbf{working specification.} Its parts deliberately interlock --- remove loops and enabling loses its mechanism, remove policing and the heat-death argument loses its filter, remove reachability and slack-flow has nothing to reshape --- so the density is \textbf{load-bearing for its function as a reference}, not incidental bloat. It is deliberately \textit{not} optimized for comfortable first reading; the cure for that is not a thinner specification (which would dissolve the interlock) but a separate, staged exposition that introduces the chain one link at a time. That companion reader-facing document exists alongside this one.
}
\end{quote}

\hr

\section{The primitive}
The thing that governs whether a new structure --- species, invention, crystal, collaboration --- comes into being is a competition between two terms with a barrier between them:

\begin{quote}
\textbf{Formation cost} (the energy / effort / improbability of bringing the structure into existence) versus \textbf{yield} --- the capacity the structure \textit{releases} once it exists and runs (its \textbf{released slack}, §10). The structure is worth forming iff yield exceeds cost --- and between the two sits a \textbf{barrier}.
\end{quote}

This is substrate-agnostic because "what did it cost to form this, and does what it releases repay the cost" is the \textit{same question} for a crystal, a species, an invention, and an institution. That sameness, with content rather than vocabulary, is the claim of Evolution by Emergence.

One flag to carry forward --- and it corrects a subtle error in earlier drafts. Yield (released slack) is a \textbf{scalar}, and that is \textit{fine}: in the cleanest physical substrate (crystals, §3) it is literally bulk free energy, one number, and in every substrate it is a conserved, budgeted quantity (§10). The scalar-ness was never the problem. What cannot be a scalar is what the slack \textit{does downstream} --- the \textbf{reshaping of which futures become reachable} --- and trying to compress \textit{that} into a single "future value" number is the teleology trap §9 forbids. So keep the split clean throughout: \textit{yield = released slack} (scalar, the cause); \textit{reachability-reshaping} = what it pays for (structural, the effect).

---

\section{Two ratios, separated --- and only one is the engine}

\begin{itemize}
    \item \textbf{Ratio A (Jevons / rebound --- deployment).} Efficiency rises 50\%, but if usage more than doubles, total consumption rises anyway. Real and important, but a fact about \textit{how released slack gets spent at the system level} --- contingent on demand and behaviour. A downstream \textit{consequence}, not the engine. Set aside for now (though §10 shows the slack it describes is the very thing that drives structure formation --- the "consequence" turns out to be the fuel).

    \item \textbf{Ratio B (formation cost vs yield --- the engine).} It costs energy to make the step \textit{exist}: 2 gallons of effort to build an improvement that saves 0.5 is a net 1.5 loss, and the improvement was a dead end. This is the cost of \textit{forming} the pathway, not running it --- the activation barrier, the "energy hurdle" from the first message. \textbf{This is the one the framework is about.}
\end{itemize}

Ratio B immediately refutes a tempting claim --- \textit{"any efficiency gain is beneficial."} No. The 2-for-0.5 case is a real efficiency gain that is \textbf{net negative}, because forming it cost more than it will ever release. Gains are beneficial \textbf{iff formation cost is repaid by yield}, and the sign can be negative. Most conceivable pathways are \textit{not} worth forming; the rare ones that are repay the barrier, and those are what we see persist. The filter is the ratio.

---

\section{The template: classical nucleation theory}
Your ratio is \textbf{already a worked physical theory} in the cleanest substrate. A crystal forming from a melt or solution faces exactly cost-versus-yield-with-a-barrier. The free energy to form a cluster of radius $r$:

\begin{verbatim}
ΔG(r)  =  (surface term, COST)      +  (volume term, YIELD)
       =  4π r² γ                   −  (4/3) π r³ Δgᵥ
\end{verbatim}

\begin{itemize}
    \item $\gamma$ = \textbf{interfacial energy} --- cost per unit area of the \textit{boundary} between the new structure and its surroundings; the cost of the bonds/seams that separate "this structure" from "not it." Positive. Scales as $r^2$.
    \item $\Delta g_v$ = \textbf{driving force} --- free energy released per unit volume of the new phase. The (scalar) yield density. Scales as $r^3$.
\end{itemize}

Because cost scales as $r^2$ and yield as $r^3$, for \textit{small} clusters the surface cost dominates and $\Delta G$ \textbf{rises}; for \textit{large} clusters the volume yield dominates and $\Delta G$ \textbf{falls}. So there is a maximum --- a \textbf{barrier} --- at the \textbf{critical radius $r*$}. Below $r*$, clusters shrink; above $r*$, they grow. A nascent structure below critical size \textit{dissolves back}; only one that clears the barrier persists. The barrier height:

\begin{verbatim}
ΔG*  ∝  γ³ / Δgᵥ²
\end{verbatim}

rises with the \textbf{cube} of the interfacial cost and falls with the \textbf{square} of the driving force. The crystal is the proof-of-concept of the whole framework: formation cost (surface) versus yield (bulk), a sharp barrier between, structure appearing \textbf{only when yield wins.}

---

\section{The mapping across substrates --- a hypothesis, not an asserted identity}
The reach is the conjecture that \textbf{many adaptive systems appear to exhibit nucleation-like barrier dynamics} --- the same three quantities ($\gamma$, $\Delta g_v$, $\Delta G*$) having \textit{recognisable analogues} across substrates. Offered as a mapping to test case by case, not a demonstrated universality. Each row could fail; the principle survives any individual failure.

\begin{center}
\begin{tabular}{>{\raggedright\arraybackslash}p{3.5cm}>{\raggedright\arraybackslash}p{3.5cm}>{\raggedright\arraybackslash}p{3.5cm}>{\raggedright\arraybackslash}p{3.5cm}}
\toprule
\textbf{Substrate} & \textbf{Interfacial cost $\gamma$ (boundary/bond cost)} & \textbf{Driving force $\Delta g_v$ (yield density)} & \textbf{What "clearing the barrier" means} \\
\midrule
\textbf{Crystal} (\textit{solved}) & surface energy of the interface & bulk free energy released by ordering & cluster exceeds critical radius, grows \\
\textbf{Species} (\textit{analogy}) & cost of crossing the fitness valley; mutational distance to a viable adaptation & persistence the new niche-occupancy buys & a viable, reproductively coherent form establishes \\
\textbf{Invention} (\textit{analogy}) & R\&D cost; effort to assemble a working configuration & capacity the working invention releases & the thing works and propagates \\
\textbf{Collaboration} (\textit{analogy}) & cost of forming trust-bonds between agents & durable capacity the bound collective releases & a stable cooperative holds against defection \\
\bottomrule
\end{tabular}
\end{center}

Only the first row is established physics. The rest are \textit{candidate} mappings, and the work (§7) is to show that for at least one, the quantities can be defined independently and predict something simpler theories do not.

---
\section{The tweakable parameters --- the three knobs nucleation theory says exist}
The barrier is \textbf{not fixed}; it is a function of parameters you can move. Nucleation theory says there are exactly three knobs, and your own examples land on them.

\begin{itemize}
    \item \textbf{Knob 1 --- lower the interfacial cost $\gamma$.} Because $\Delta G* \propto \gamma^3$, this is the \textit{steepest} lever. Your examples --- \textit{"less hate," "less suspicion," "more honesty"} --- are exactly this in the collaboration substrate: the interfacial cost between agents \textit{is} the cost of forming trust-bonds; hate and suspicion raise $\gamma$, honesty and shared norms lower it. The cube law makes this the most powerful intervention available --- reducing the cost of trust cuts the barrier \textit{cubically}.

    \item \textbf{Knob 2 --- raise the driving force $\Delta g_v$ (supersaturation).} Push further from equilibrium --- more latent, unexpressed potential --- and the barrier falls as $1/\Delta g_v^2$. \textit{"The world is soooo ready"} is a claim of high supersaturation. But high supersaturation \textbf{lowers} the barrier, it does not \textbf{abolish} it: a supersaturated solution with no seed sits metastable indefinitely. Readiness is necessary, not sufficient.

    \item \textbf{Knob 3 --- provide a template (heterogeneous nucleation).} A foreign surface or seed \textbf{pays part of the interfacial cost for you} --- why dust nucleates raindrops, why seeding works. \textbf{Your oak proposal is heterogeneous nucleation:} bringing oak species into proximity to cross-breed both raises the supersaturation of variation (Knob 2) and provides templates --- existing viable genomes a new climate-resilient configuration can assemble \textit{onto} rather than from nothing (Knob 3).
\end{itemize}

\textbf{A fact that strengthens the oak case:} oaks (\textit{Quercus}) are the textbook \textbf{syngameon} --- distinct species that stay ecologically separate yet exchange genes across boundaries through hybridisation. The oak syngameon exhibits dynamics closely resembling natural heterogeneous nucleation from a shared interfertile gene pool.

---
\section{The brakes --- lowering the barrier is not free, and two examples have a known failure mode}
The barrier is also a \textbf{filter.} Lower it too far and you don't get \textit{better} structures faster --- you nucleate \textbf{junk.}

\begin{itemize}
    \item \textbf{A sweet spot for $\gamma$, not a floor at zero.} Suspicion is not \textit{only} interfacial cost --- it is also a \textbf{defector-detection mechanism.} Drop it to zero and collaborations nucleate readily, then get \textbf{eaten by free-riders}. $\gamma$ cannot go to zero without removing what makes the structure durable.

    \item \textbf{Too much supersaturation nucleates defective structures.} Push $\Delta g_v$ too high and you get \textit{many small unstable} nuclei rather than a few good ones. In the oak case this is \textbf{outbreeding depression:} cross too freely and you swamp locally-adapted gene complexes.

    \item \textbf{The template imposes its own configuration.} A bad seed nucleates a bad structure efficiently (path dependence, lock-in).
\end{itemize}

---
\section{What to do next --- the falsification challenge}
Show, in \textbf{at least one non-physical domain}, that formation cost, driving force, barrier height, and maintenance/policing cost can be defined \textbf{independently of each other} and \textbf{independently of the outcome being predicted}, and that the framework then predicts something a simpler existing theory would not.

---
\section{The level-jump: why stacking structure around \textit{agents} is categorically harder}
Structure-formation is \textbf{recursive}. The cost term acquires a component that exists only when the components are agents:

\begin{verbatim}
γ(level-jump)  =  cost of forming the bonds  +  cost of POLICING the components against defection
\end{verbatim}

That second term \textbf{does not exist when components are inert.} If policing enters $\gamma$ additively, then because $\Delta G* \propto \gamma^3$, a high policing cost makes the barrier brutal and the transition rare.

---
\section{What the yield \textit{does} --- the structural reshaping of reachability}
\textbf{Move one:} Yield is released slack (§10) --- the budgeted, conserved, \textit{scalar} capacity a structure frees.

\textbf{Move two:} The released slack \textbf{reshapes the landscape/graph itself}, $(V, E) \rightarrow (V, E')$: the released slack lowers some barriers, and \textit{which futures become reachable} changes.

The enabling is structural, not scalar --- and that is what dissolves the teleology.

---
\section{The mechanism --- loops and slack-flow}
\textbf{Structure is process --- maintenance is constitutive.} The structure \textit{is} the maintenance, the way a flame \textit{is} the burning.

\textbf{Slack-flow is the enabling mechanism --- and it is local.} Cheapening a loop \textbf{releases capacity now}, and that capacity lowers the barrier to \textit{other} loops.

\textbf{Three moments in one flow:} Formation pays to \textit{ignite} the loop. Maintenance pays to \textit{keep it running}. Enabling is the \textit{slack released} by running it more efficiently.

---
\section{The network layer --- two gradients that must not be merged}
\textbf{Gradient 1 --- descent.} A network moves down an \textbf{error surface}.

\textbf{Gradient 2 --- nucleation.} Noise-assisted hop \textbf{is the thermal-fluctuation barrier-crossing of nucleation.}

---
\section{The flame --- what "prolonging the existence of the network" actually means}
A **rock** persists \textit{for free}. A **flame** persists \textit{only by continuous payment} --- fuel, oxygen, every instant. \textbf{A network is a flame, not a rock.}

---
\section{What the framework is about --- structure \textit{and} flow}
The framework's non-triviality lives entirely in the places where flow is \textit{interrupted}, not where it flows. Every load-bearing element is a \textit{resistance} to flow, a boundary, a held structure.

\textbf{A structure is a loop that pays to persist \textit{against} the flow it is made of.}

---
\section{What \textit{kind} of claim this is --- an accounting identity plus a dynamics}
The framework's spine is an \textbf{accounting identity}: every persistent structure must balance a formation-and-maintenance ledger.

\textbf{The falsifiable, non-trivial content is entirely in the dynamics hanging off the identity.}

\hr

\textit{Stay suspicious of the warm click of recognition. The disciplines that keep this real: (1) When lowering a barrier, name which part of $\gamma$ is friction and which is filter. (2) When calling something valuable for what it enables, force the claim into the reshaping-of-reachability form. (3) Never optimize reachability or viability directly. (4) Keep the gradient present-tense. (5) Lean on the shared hard mathematics. (6) Do not mistake the framework's substrate for its subject.}

\end{document}